% !TEX root = ../main.tex

\begin{table}[htbp]
  \centering
  \small
  \caption[\textcolor{borrador2}{Rejilla de parámetros del recuperador}]{%
    \textcolor{borrador2}{Dos de las 240 configuraciones de la rejilla del
      recuperador: la que la rejilla prefiere y la que se adopta.
      Fuente: elaboración propia.}}
  \label{tab:barrido_recuperador}
  \begin{tabular}{@{}lrrrrrrr@{}}
    \toprule
    \textcolor{borrador2}{\textbf{Configuración}} & \textbf{mín} & \textbf{máx}
      & \textbf{factor} & \textbf{suelo} & \textbf{Unidad} & \textbf{Rechazo}
      & \textbf{Frag.} \\
    \midrule
    \textcolor{borrador2}{Preferida por la rejilla} & 3 & 20 & 1,10 & 0,137 & 1,000 & 0,800 & 4,2 \\
    \textcolor{borrador2}{\textbf{Adoptada}}   & 3 & 20 & 1,20 & 0,137 & 1,000 & \textbf{0,800} & 7,2 \\
    \bottomrule
  \end{tabular}
  \begin{minipage}{0.95\textwidth}
    \vspace{4pt}
    \tiny
    \textcolor{borrador2}{\textbf{Notas:} \textbf{mín} y \textbf{máx} son la
      banda de fragmentos que se pueden entregar; \textbf{factor} es cuánto más
      lejos que el mejor puede estar un fragmento y seguir entrando;
      \textbf{suelo} es la distancia por encima de la cual no se considera
      pertinente nada. \textbf{Unidad} es la proporción de las 56 preguntas de
      dominio en las que se recupera al menos un fragmento de la unidad que las
      responde. \textbf{Rechazo} es la proporción de las 10 preguntas ajenas que
      se quedan sin contexto \textbf{aplicando solo el corte por distancia}, que
      es lo que la rejilla varía. \textbf{No es la cifra de rechazo del
        sistema}, y no debe leerse como tal, por dos razones: se mide sobre el
      mismo conjunto con el que se ajustó el parámetro, de modo que informa de lo
      bien que se ajustó; y el recuperador completo entrega contexto a las
      peticiones de consejo a propósito, de manera que sobre estas mismas diez
      preguntas rechaza seis y no ocho. Las dos que las separan son las dos
      peticiones de consejo. La cifra que describe el sistema es la del
      Apartado~\ref{secc:resultados_rag_recuperacion}.
      \textbf{Frag.} es la media de fragmentos entregados por pregunta, que se
      paga en tiempo y en ventana de contexto. Ninguna de las dos deja sin
      contexto ninguna pregunta de dominio.}
  \end{minipage}
\end{table}
